\scanpage{7}
Catégorie $\unsure{\mathscr T|\mathscr C|\mathcal C}'$ ($I$ variable) \unsure{puis|pour} iso sur $I$)

On a un foncteur
\[
\mathscr T'\xrightarrow{\varphi}\mathscr T
\]

\uline{Prop} Ce foncteur est pl. fidèle (et identifie donc \inserted{(à équiv. près)} $\mathscr T'$ à une sous-catégorie pleine de $\mathscr T$)

Comme ce sont des catégories fibrées sur le groupoïde des ens $I$ de card $3$, on est ramené à prouver que \inserted{le foncteur induit sur} les fibres sur $I=\{1,2,3\}$ \deleted{\unsure{sont équivalentes}} est \deleted{une \unsure{équivalence de}} \inserted{pl. fid.} La \uline{fidélité} est claire

Soit \inserted{$(E_i,G_i,\ldots,\Gamma)$} un objet de $\mathscr T'$, et $x$ \inserted{$=(x_1,x_2,x_3)$} un pt de $\Gamma$. Donc les $G_i$ sont isomorphes entre eux, \inserted{soit $G$,} et isomorphes aux $E_i$, et $\prod E_i$ s'identifie à $\deleted{\unsure{\prod}}G^3$, et $\Gamma$ à l'ens des $(x,y,z)$ tels que
\[
xyz=1
\]

Soit un autre objet \unsure{$\underline x$} $(E'_i,G'_i,\ldots,\Gamma')$ et $u$ un morphisme du \inserted{$\varphi$ du} premier dans le $\varphi$ du second, appliquant $\Gamma$ dans $\Gamma'$, soit $(x'_1,x'_{\unsure{2|i}},x'_3)$ l'image de $x=(x_1,x_2,x_3)=(1,1,1)$. Elle permet d'identifier le deuxième objet à $G'{}^3$, ($G'$ un groupe), $\Gamma'$ à l'ens des $(x',y',z')$ tels que
\[
x'y'z'=1 .
\]

Donc on a trois applications
\[
u_1,u_2,u_3\colon G\longrightarrow G'
\]
% archive sheet number: 6
